\begin{prob}
    Determine the worst case time complexity of each of the recursive algorithms below.
    In each case, state the recurrence relation describing the runtime. Solve the recurrence
    relation, either by unrolling it or showing that it is the same as a recurrence we have
    encountered in lecture.

    \begin{subprobset}
        \begin{subprob}
            \inputminted{python}{\thisdir/include/part_a.py}

            \begin{soln}
                Note that the slicing in \python{numbers[:mid_left]} (and so on) takes linear time
                in the length of the array, so the recurrence relation is:
                \[
                    T(n) = \begin{cases} 
                        \Theta(n) + 3 T(n/3) & n > 1 \\
                        \Theta(1) & n \leq 1
                     \end{cases}
                \]
                This is the same as the  recurrence relation as for 3-way mergesort, and the solution
                is $\Theta(n \log n)$.
            \end{soln}
        \end{subprob}


        \begin{subprob}

            \inputminted{python}{\thisdir/include/part_b.py}

            \begin{soln}
                No slicing is being done, so a constant amount of time is required outside
                of the recursive calls. Therefore the recurrence relation is:
                \[
                    T(n) = \begin{cases} 
                        \Theta(1) & \text{stop} \leq \text{start or stop} - \text{start} = 1 \\
                        2T(n/2) + \Theta(1) & \text{otherwise}
                     \end{cases}
                \]

                We have not seen this recurrence before, so we must solve it by unrolling. 
                For simplicity, assume $T(n)= 2T(n/2) + 1$. 
                If we perform unrolling $k$ times, we will get the following: 
                \[ T(n) = 2^k T(n/2^k) + 1 + 2 + 4 + \cdots + 2^{k-1}.
                \] 
                This unrolling stops when we have $n / 2^k = 1$, meaning $2^k = n$ and $k = \log_2 n$. 
                Furthermore, note that $1 + 2 + 4 + 8 + \cdots + 2^{k-1} = 2^k - 1$ (geometric sum). 
                We thus have 
                \[ T(n) = n T(1) + 2^k  - 1 = n \cdot \Theta(1) + n - 1 = \Theta(n). \] 
                It will have $\Theta(n)$ as its solution.

                Notice that this recurrence is similar to the recurrence for
                binary search, which was $T(n/2) + \Theta(1)$, and whose solution was
                $\Theta(\log n)$. The difference between the two is the factor
                of 2 on the $T(n/2)$; binary search does not have this because it only
                makes one recursive call. This difference has a big effect: it reduces
                the time complexity from $\Theta(n)$ all the way down to $\Theta(\log n)$.

            \end{soln}
        \end{subprob}

        \begin{subprob}
            In this problem, remember that \python{//} performs \emph{flooring division},
            so the result is always an integer. For example, \python{1//2} is zero.
            \python{random.randint(a,b)} returns a random integer in $[a,b)$ in constant time. 
            Assume \python{isEven()} takes in an integer and returns whether it is even
            in constant time. Note that you are asked to determine the {\bf worst-case} time complexity of the following algorithm. 

            \inputminted{python}{\thisdir/include/part_c.py}

            \begin{soln}
                The work outside of the recursive calls takes linear time. (Note that while there are random numbers generated, the time complexity for outside the recursive calls are deterministic.) 
                Depending on whether $x$ is even or not, the algorithm will make {\bf one} recursive call on problem of size $n/2$. Therefore the recurrence is: 
                \[
                T(n) = \begin{cases} 
                    T(n/2) + \Theta(n) & n > 0 \\
                    \Theta(1) & n = 0
                 \end{cases}
                \]
                We can solve this by unrolling it. In particular, for simplicity, assume $T(n) = T(n/2) + cn$. 
                Applying unrolling $k$ times, we have have the following: 
                \[ T(n) = T(\frac{n}{2^{k}}) + c \frac{n}{2^{k-1}} + c\frac{n}{2^{k-2}} + \cdots + c\frac{n}{2} + cn . \]
                The unrolling stops when we have $n / 2^k = 1$, meaning that $2^k = n$ and $k=\log_2 n$. Furthermore,  
                \[ n + n/2 + \cdots n/ 2^{k-1} = n (1 + 1/2 + 1/4 + \cdots + 1/2^{k-1}) = n \cdot \Theta(1) = \Theta(n).\]
                It then follows that for $k = \log_2 n$, we have that 
                \[T(n) = T(\frac{n}{2^{k}}) + c \frac{n}{2^{k-1}} + c\frac{n}{2^{k-2}} + \cdots + c\frac{n}{2} + cn = \Theta(1) + \Theta(n) = \Theta(n).\]
                Therefore $T(n) = \Theta(n)$. 
              
            \end{soln}
        \end{subprob}

    \end{subprobset}

\end{prob}
